\documentclass[10pt]{article}

% ---------------------------------------------------------------------------
% Packages
% ---------------------------------------------------------------------------
\usepackage[T1]{fontenc}
\usepackage[utf8]{inputenc}
\usepackage{lmodern}
\usepackage[margin=1in]{geometry}
\usepackage{amsmath}
\usepackage{graphicx}
\usepackage{hyperref}
\usepackage{xcolor}
\usepackage{listings}
\usepackage{booktabs}
\usepackage{caption}
\usepackage{subcaption}
\usepackage{microtype}
\usepackage{natbib}
\usepackage{authblk}
\usepackage{abstract}
\usepackage{titlesec}
\usepackage{enumitem}
\usepackage{float}

% ---------------------------------------------------------------------------
% Listing style (Python / shell)
% ---------------------------------------------------------------------------
\definecolor{codegray}{rgb}{0.95,0.95,0.95}
\definecolor{codeblue}{rgb}{0.00,0.00,0.60}
\definecolor{codegreen}{rgb}{0.00,0.50,0.00}
\definecolor{codemauve}{rgb}{0.58,0.00,0.82}

\lstdefinestyle{python}{
  backgroundcolor=\color{codegray},
  basicstyle=\ttfamily\footnotesize,
  breakatwhitespace=false,
  breaklines=true,
  captionpos=b,
  commentstyle=\color{codegreen},
  frame=single,
  keepspaces=true,
  keywordstyle=\color{codeblue},
  language=Python,
  numbers=none,
  showspaces=false,
  showstringspaces=false,
  stringstyle=\color{codemauve},
  tabsize=2,
  framerule=0.4pt,
  xleftmargin=4pt,
  xrightmargin=4pt,
}

\lstdefinestyle{bash}{
  backgroundcolor=\color{codegray},
  basicstyle=\ttfamily\footnotesize,
  breaklines=true,
  frame=single,
  numbers=none,
  framerule=0.4pt,
  xleftmargin=4pt,
  xrightmargin=4pt,
}

% ---------------------------------------------------------------------------
% Hyperref config
% ---------------------------------------------------------------------------
\hypersetup{
  colorlinks=true,
  linkcolor=codeblue,
  citecolor=codegreen,
  urlcolor=codeblue,
}

% ---------------------------------------------------------------------------
% Metadata
% ---------------------------------------------------------------------------
\title{\textbf{MetaKG: A Unified Knowledge Graph System\\for Metabolic Pathway Analysis}}

\author[1]{Eric G. Suchanek, PhD}
\affil[1]{Flux Frontiers, \href{https://github.com/Flux-Frontiers}{github.com/Flux-Frontiers}}

\date{February 2026}

% ---------------------------------------------------------------------------
\begin{document}
% ---------------------------------------------------------------------------

\maketitle

% ---------------------------------------------------------------------------
\begin{abstract}
\noindent
We present \textsc{MetaKG}, a local-first metabolic knowledge graph whose
dual-layer architecture unifies structural precision and semantic expressivity
in a single system: SQLite encodes the full property graph of compounds,
reactions, enzymes, and pathways with stoichiometric and regulatory
relationships, while LanceDB maintains dense vector embeddings for
natural-language semantic search over the same entities. This design delivers
four orthogonal query modalities---natural-language pathway discovery,
structural neighbourhood traversal, shortest-path search between metabolites,
and stoichiometric detail assembly---without forcing the user to choose between
relational precision and semantic exploration. A post-build name-enrichment
pipeline replaces bare KEGG accessions with human-readable names sourced from
KEGG bulk name lists, making all query results immediately interpretable.
Beyond querying, \textsc{MetaKG} integrates three simulation modalities
directly into the knowledge graph: Flux Balance Analysis for steady-state flux
optimisation, kinetic ODE integration with Michaelis-Menten rate laws (seeded
from BRENDA and SABIO-RK) using an implicit stiff solver (BDF), and what-if
perturbation analysis for enzyme knockouts and substrate overrides. Interactive
exploration is supported by a 2D Streamlit graph browser and a 3D PyVista
visualiser with Fibonacci-annulus and layer-stratified layouts (both under
active development). All capabilities are exposed through a unified Python API,
command-line interface, and Model~Context~Protocol (MCP) server exposing nine
typed tools, making \textsc{MetaKG} a first-class data source for AI reasoning
systems. We demonstrate the system on the complete \textit{Homo sapiens}
metabolome---all 369 KEGG pathways ingested as a single queryable graph
(17,050 nodes, 11,286 edges, 9,427 enzymes, 20,151 indexed embeddings)---with
structural queries completing in 10--50\,ms, semantic searches in
100--500\,ms, and kinetic simulations in under 400\,ms on commodity
hardware.\footnote{All timings measured on a 2024 MacBook Pro (Apple M3
Max, 36\,GB RAM, 1\,TB SSD).}
\end{abstract}

\textbf{Keywords:} metabolic pathways, knowledge graph, dual-layer architecture,
semantic search, local-first systems, AI-accessible data, multi-format integration,
bioinformatics

% ---------------------------------------------------------------------------
\section{Introduction}
% ---------------------------------------------------------------------------

\subsection{The Problem: Query Architecture Limits Exploratory Analysis}

Reconstructing metabolic networks means integrating data from diverse sources:
KEGG provides human-curated pathway maps in KGML~\citep{kanehisa2021kegg};
Reactome and model-organism databases export in
BioPAX Level~3~\citep{demir2010biopax}; constraint-based modelling tools
produce SBML~\citep{keating2020sbml}; and institutional data often lives only
in spreadsheets. Handling multiple formats is a technical necessity, but it is
not the core scientific problem. The deeper issue is \emph{query capability}.

Existing systems force a choice between two incompatible paradigms:

\begin{description}[leftmargin=*,itemsep=2pt]
  \item[Structural precision.] Systems like KEGG, BioCyc, and Reactome support
    exact structural queries---``retrieve all products of this enzyme'' or
    ``find the shortest path between two compounds''---through relational data
    models or graph databases. Yet they lack semantic search: a biologist
    looking for ``fatty-acid oxidation'' gets keyword matches, not conceptually
    similar pathways. They also require web access and resist programmatic
    querying at scale.

  \item[Semantic expressivity.] Vector databases and embeddings handle
    semantic similarity well---``find pathways similar to X''---and cope with
    synonymy and domain terminology. But they sacrifice property graph
    structure; stoichiometric detail, cofactor roles, and regulatory logic
    disappear. Semantic-only systems are exploration tools, not precision
    instruments.
\end{description}

No existing system unites structural precision, semantic expressivity, and
programmatic accessibility in a single local tool. KEGG and PathBank are
web-only. MetaNetX reconciles identifiers but offers no queryable graph.
Neo4j-based systems handle complex structural queries but require server
deployment and do not address parsing. Vector databases excel at search but
discard relational structure. Put simply, no platform lets an analyst start
with a semantic query (``what pathways relate to glucose metabolism?''), drill
into structural detail (``what is the precise stoichiometry of this
reaction?''), and assemble the results---all within a single local,
reproducible, programmable system.

\subsection{The Solution: Dual-Layer Query Architecture}

\textsc{MetaKG} resolves this tension with a dual-layer local knowledge graph
that eliminates the forced choice between structural precision and semantic
expressivity. The key idea is to separate query problems by modality:

\begin{description}[leftmargin=*,itemsep=2pt]
\item[Structural layer.] SQLite stores the property graph (compounds,
  reactions, pathways, and their relationships). Structural queries run as
  efficient SQL joins: neighbourhood traversal, shortest-path BFS,
  stoichiometric assembly.
\item[Semantic layer.] LanceDB maintains a vector index over node descriptions
  using sentence-transformer embeddings. Semantic queries retrieve results by
  cosine similarity: natural-language pathway search, compound similarity,
  synonym resolution.
\end{description}

Together, these layers give analysts access to four query modalities on the
same graph, without committing to a single paradigm. A researcher can begin
with semantic exploration and then drill into structural detail, all in one
interface:

\begin{enumerate}[leftmargin=*, itemsep=2pt]
\item \textbf{Semantic pathway discovery} --- ``Find pathways related to
  fatty-acid beta-oxidation'' (vector similarity search; handles synonyms,
  abbreviations, and domain terminology)
\item \textbf{Structural neighbourhood traversal} --- ``Find all compounds
  that are products of pyruvate carboxylase'' (SQL joins)
\item \textbf{Shortest-path search} --- ``What is the minimal metabolic route
  from glucose to acetyl-CoA?'' (BFS over the compound--reaction bipartite
  graph)
\item \textbf{Stoichiometric detail assembly} --- ``Retrieve all substrates,
  products, enzymes, inhibitors, and cofactors for reaction R00200''
  (multi-table JOIN with JSON unpacking)
\end{enumerate}

Multi-format parsing (KGML, SBML, BioPAX, CSV) is the enabling
infrastructure. It lets users ingest data from any source and merge it through
stable, deterministic identifiers. But parsing is table stakes. The innovation
is the dual-layer query engine: it frees researchers from choosing between
semantic expressivity and structural precision.

All four query modalities are exposed through a unified Python API,
command-line interface, and Model~Context~Protocol (MCP) server. The system
runs entirely locally---no network, no database server, no external services
after setup. This is a deliberate design choice. Unlike live database mirrors,
\textsc{MetaKG} treats the knowledge graph as a reproducible,
version-controlled snapshot. Users rebuild when input files change, enabling
reproducible analysis, offline workflows, and integration into research
pipelines. The system scales to organism-scale pathway corpora: the complete
\textit{Homo sapiens} metabolome (369 KEGG pathways, 17,050 nodes, 11,286
edges, 9,427 enzymes) builds in 30--60~seconds, and queries complete in
milliseconds to seconds.

\subsection{Design Goals}

\textsc{MetaKG} is built around three core principles, in priority order:

\begin{enumerate}[leftmargin=*, itemsep=2pt]
\item \textbf{Dual-layer query architecture.} Combine structural and semantic
  queries in a single system. Do not force users to choose between precise
  graph traversal and exploratory semantic search; both modalities should work
  seamlessly on the same data through the same API.

\item \textbf{Format-agnostic, deterministic merging.} Accept pathway data in
  any of four formats (KGML, SBML, BioPAX, CSV) and produce a unified graph
  with stable, deterministic identifiers. Enable reproducible builds and
  incremental updates without external reconciliation services.

\item \textbf{Zero-friction local deployment.} Require no database servers,
  external services, or network connections after setup. The entire
  stack---parsing, storage, indexing, querying, visualisation, MCP
  server---fits in a single Python library, runnable on a laptop or in batch
  workflows.
\end{enumerate}

The remainder of this paper is organised as follows. Section~\ref{sec:design}
describes the data model and system architecture. Section~\ref{sec:parsers}
covers the format-specific parsers. Section~\ref{sec:enrich} describes the
name-enrichment pipeline. Section~\ref{sec:storage} details the storage and
indexing layers. Section~\ref{sec:query} describes the query API.
Section~\ref{sec:viz} covers the visualisation components.
Section~\ref{sec:mcp} describes the MCP server. Section~\ref{sec:usage}
walks through a complete worked example. Section~\ref{sec:discussion}
discusses limitations and future directions, and Section~\ref{sec:conclusion}
concludes.

% ---------------------------------------------------------------------------
\section{Data Model and Architecture}
\label{sec:design}
% ---------------------------------------------------------------------------

\subsection{Property Graph Schema}

\textsc{MetaKG} represents metabolism as a directed property graph
$G = (V, E)$. Vertices $V$ are \emph{entities} of four kinds:

\begin{description}[leftmargin=*, itemsep=2pt]
  \item[\texttt{compound}] A small-molecule metabolite. May carry a molecular
    formula, net formal charge, and cross-references to external databases
    (ChEBI, HMDB, PubChem, InChI).
  \item[\texttt{reaction}] A biochemical transformation. Carries stoichiometry
    encoded as a JSON object listing substrates and products with their
    coefficients, a reversibility flag, and cross-references to KEGG and Rhea.
  \item[\texttt{enzyme}] A protein catalyst. Carries an EC number and
    cross-references to UniProt and NCBI Gene.
  \item[\texttt{pathway}] An ordered or thematic collection of reactions,
    typically corresponding to one named pathway in a source database.
\end{description}

Edges $E$ carry one of seven relation types (Table~\ref{tab:relations}).
All edges are directed and may include an evidence blob encoded as JSON,
which parsers use to record stoichiometric coefficients, compartment labels,
and source-specific annotations.

\begin{table}[H]
  \centering
  \caption{Edge relation types in the MetaKG graph schema.}
  \label{tab:relations}
  \small
  \begin{tabular}{@{}lll@{}}
    \toprule
    Relation & Source kind & Target kind \\
    \midrule
    \texttt{SUBSTRATE\_OF}  & compound  & reaction  \\
    \texttt{PRODUCT\_OF}    & reaction  & compound  \\
    \texttt{CATALYZES}      & enzyme    & reaction  \\
    \texttt{INHIBITS}       & compound  & reaction  \\
    \texttt{ACTIVATES}      & compound  & reaction  \\
    \texttt{CONTAINS}       & pathway   & reaction  \\
    \texttt{XREF}           & any       & any       \\
    \bottomrule
  \end{tabular}
\end{table}

\subsection{Stable Identifier Scheme}

Identifier reconciliation is one of the persistent difficulties in biological
data integration. \textsc{MetaKG} assigns each node a stable, URI-style
string identifier of the form:

\begin{center}
\texttt{<prefix>:<namespace>:<external-id>}
\end{center}

\noindent where the prefix encodes the node kind (\texttt{cpd}, \texttt{rxn},
\texttt{enz}, \texttt{pwy}) and the namespace identifies the source database
(\texttt{kegg}, \texttt{chebi}, \texttt{uniprot}, \texttt{ec}, etc.):

\begin{lstlisting}[style=python, caption={Examples of stable node identifiers.}]
cpd:kegg:C00022       # Pyruvate (KEGG)
rxn:kegg:R00200       # Pyruvate decarboxylation
enz:ec:1.2.4.1        # Pyruvate dehydrogenase (EC)
pwy:kegg:hsa00010     # Glycolysis / Gluconeogenesis
\end{lstlisting}

When an entity appears in a file without an external database identifier, the
system constructs a synthetic identifier by hashing the lowercased display name
with SHA-1 and retaining the first eight hexadecimal digits:

\begin{center}
\texttt{cpd:syn:a4f2b8c1}
\end{center}

Because the hash operates on the normalised name, identifiers are
deterministic across independent parser runs on the same input. This makes
incremental rebuilds possible without producing duplicate nodes. When two
source files refer to the same KEGG or ChEBI entry, the graph merges their
nodes by ID before writing to SQLite, so each logical entity appears exactly
once regardless of how many files contributed to it.

\subsection{System Architecture and the Dual-Layer Query Engine}

Figure~\ref{fig:arch} shows the overall pipeline. The \texttt{MetaKG}
orchestrator class owns the full sequence from raw files to query results,
coordinating three subsystems:

\begin{enumerate}[leftmargin=*]
  \item \textbf{MetabolicGraph} --- handles file discovery and parser
    dispatch. Outputs normalised \texttt{MetaNode} and \texttt{MetaEdge}
    objects that are independent of source format.
  \item \textbf{MetaStore} --- the SQLite persistence layer for structural
    graph queries (neighbourhood traversal, shortest-path BFS, stoichiometric
    assembly) via efficient SQL joins. Provides ACID guarantees with no
    external dependencies.
  \item \textbf{MetaIndex} --- the LanceDB vector index layer for semantic
    queries (natural-language pathway search, compound similarity retrieval).
    Uses pre-trained sentence-transformer embeddings to handle synonymy,
    abbreviations, and domain terminology.
\end{enumerate}

The dual-layer design is the core architectural contribution. By routing
structural queries (naturally SQL problems) to SQLite and semantic queries
(naturally vector problems) to LanceDB, \textsc{MetaKG} avoids forcing a
choice between relational precision and semantic expressivity. A user can
traverse the graph by stoichiometry using SQL, then search for semantically
similar pathways using embeddings---all within the same interface.

Both the store and the index initialise lazily: a caller that builds only the
SQLite database (e.g.\ with \texttt{--no-index}) never loads the embedding
model. This supports lightweight deployments where only structural queries are
needed, avoiding the 100~MB embedding model download unless semantic search is
actually used.

\begin{figure}[H]
  \centering
  \begin{lstlisting}[style=bash, frame=single, caption={High-level data flow.}]
Pathway files (KGML / SBML / BioPAX / CSV)
        |
  MetabolicGraph
  (file discovery + parser dispatch)
        |
  MetaNode / MetaEdge objects
  (merged by stable ID)
        |
  MetaStore  ----> SQLite
  (write + xref index)
        |
  MetaIndex  ----> LanceDB
  (sentence-transformer embeddings)
        |
  Name Enricher (optional post-build pass)
  Phase 1: enzyme labels from CATALYZES edges
  Phase 2: KEGG TSV -> human-readable names
        |
  Query API + CLI + MCP server
  \end{lstlisting}
  \caption{MetaKG data pipeline. Format-specific parsers produce a stream of
  normalised \texttt{MetaNode} and \texttt{MetaEdge} objects that are merged
  by stable identifier, persisted to SQLite, and indexed as dense vectors in
  LanceDB. An optional name-enrichment pass replaces bare KEGG accessions with
  human-readable names.}
  \label{fig:arch}
\end{figure}

% ---------------------------------------------------------------------------
\section{Format-Specific Parsers}
\label{sec:parsers}
% ---------------------------------------------------------------------------

\textsc{MetaKG} ingests metabolic pathway data from four standard formats:
KEGG Markup Language (KGML), Systems Biology Markup Language (SBML), Biological
Pathway Exchange (BioPAX), and plain tabular files (CSV/TSV). All parsers
conform to an abstract interface: they detect the file type, parse it, and
produce normalised \texttt{MetaNode} and \texttt{MetaEdge} objects. Parsing is
stateless and deterministic---the same input file always yields the same
output.

Format detection relies on content analysis (the root XML element) rather than
file extension, making the system robust to files served without standard
extensions. The parser dispatcher examines each file, selects the appropriate
handler, and produces a unified stream of nodes and edges. These are then
merged by stable identifier in the MetaStore layer. Detailed parser
specifications for each format appear in Appendix~\ref{sec:parsers-detail}.

% ---------------------------------------------------------------------------
\section{Name Enrichment}
\label{sec:enrich}
% ---------------------------------------------------------------------------

KGML-sourced graphs initially carry bare KEGG accessions as node names
(e.g.\ \texttt{C00031}, \texttt{R00710}). While these identifiers are
unambiguous, they are not human-readable and impede exploratory analysis.
\textsc{MetaKG} provides an optional post-build enrichment pipeline that
replaces accessions with canonical names and stores them directly in
\texttt{meta\_nodes.name}, making them available everywhere---CLI, MCP,
Streamlit, and the Python API.

\subsection{Two-Phase Enrichment}

Enrichment proceeds in two phases, both idempotent and safe to re-run:

\begin{description}[leftmargin=*,itemsep=3pt]
  \item[Phase~1 (no network, always runs).] For every reaction node still
    carrying a bare accession name, the enricher looks up all
    \texttt{CATALYZES} edges pointing to it and joins the catalysing enzyme
    gene symbols with a slash separator:
    \begin{center}
      \texttt{R00710} $\to$ \texttt{ADH1A / ADH1B / ADH1C}
    \end{center}
    This phase requires no external data and runs immediately after build.

  \item[Phase~2 (requires downloaded KEGG name lists).] KEGG publishes bulk
    name lists for compounds ($\sim$19,500 entries) and reactions
    ($\sim$12,400 entries). A provided download script fetches these files
    once; subsequent enrichment runs read them locally:
    \begin{lstlisting}[style=bash]
python scripts/download_kegg_names.py
metakg-enrich --db .metakg/meta.sqlite
    \end{lstlisting}
    Phase~2 updates compound names (e.g.\ \texttt{C00031} $\to$
    \texttt{D-Glucose}) and reaction names (e.g.\ \texttt{R00710} $\to$
    \texttt{Acetaldehyde:NAD\textsuperscript{+} oxidoreductase}), overriding
    Phase~1 enzyme labels where canonical KEGG names are available.
\end{description}

Both phases can be combined with the build step via the \texttt{--enrich}
flag:

\begin{lstlisting}[style=bash]
metakg-build --data ./pathways --enrich
\end{lstlisting}

After enrichment, all query results---including MCP tool responses and
Streamlit displays---show human-readable names rather than bare accessions.
The enrichment is stored in the database, so it persists across sessions
without re-running.

% ---------------------------------------------------------------------------
\section{Storage and Indexing}
\label{sec:storage}
% ---------------------------------------------------------------------------

The dual-layer storage architecture is what gives MetaKG its query
flexibility. Rather than routing all queries through a single backend (as
relational databases do) or abandoning structure entirely (as vector-only
systems do), it maintains two complementary stores, each optimised for a
different class of query:

\begin{description}[leftmargin=*,itemsep=3pt]
\item[SQLite] handles efficient relational queries over graph structure.
  Compound--reaction--pathway relationships map naturally to SQL joins.
  Shortest-path searches use BFS over edges. Stoichiometric detail
  (substrates, products, coefficients) is stored as decoded JSON columns.

\item[LanceDB + embeddings] handles semantic search over node descriptions.
  Users query in natural language (``fatty-acid oxidation''), and the system
  returns semantically similar pathways via vector similarity. This matters
  most for exploratory analysis where the user does not know exact identifiers.
\end{description}

The two layers are not redundant---they solve fundamentally different
problems. A researcher might start with semantic discovery (``which pathways
involve energy metabolism?''), then drill into structural detail (``what is the
precise stoichiometry of ATP synthesis?''), all without leaving the interface.

\subsection{SQLite Layer}

Parsed nodes and edges are written to a SQLite database through the
\texttt{MetaStore} class. The schema uses three tables:

\begin{description}[leftmargin=*,itemsep=2pt]
  \item[\texttt{meta\_nodes}] One row per node. Columns correspond to the
    fields of \texttt{MetaNode}: \texttt{id}, \texttt{kind}, \texttt{name},
    \texttt{description}, \texttt{formula}, \texttt{charge},
    \texttt{ec\_number}, \texttt{stoichiometry} (JSON), \texttt{xrefs}
    (JSON), \texttt{source\_format}, \texttt{source\_file}. Indexed on
    \texttt{kind} and \texttt{name}.
  \item[\texttt{meta\_edges}] One row per directed edge: \texttt{src},
    \texttt{rel}, \texttt{dst}, \texttt{evidence} (JSON). Indexed on
    \texttt{src}, \texttt{dst}, and \texttt{rel}.
  \item[\texttt{xref\_index}] A materialised inverse mapping from each
    external identifier to the corresponding internal node ID, built after
    all nodes have been written. This enables look-up by KEGG compound ID,
    ChEBI accession, UniProt accession, or any other cross-reference stored
    in the \texttt{xrefs} JSON blob.
\end{description}

SQLite runs with write-ahead logging (\texttt{WAL} journal mode) and
\texttt{NORMAL} synchronisation; these pragmas deliver throughput close to an
in-memory database while preserving crash safety for workloads that write once
and read many times.

\subsection{Semantic Index}

Structural queries alone are insufficient for exploratory use cases where the
user does not know an exact identifier. \textsc{MetaKG} therefore maintains a
vector index over node descriptions using LanceDB~\citep{lancedb2023} as the
approximate nearest-neighbour engine and the \texttt{all-MiniLM-L6-v2}
sentence-transformer model~\citep{reimers2019sbert} as the default encoder.
This model produces 384-dimensional embeddings and occupies approximately
100~MB of disk space on first download.

Each node is serialised to an embedding string that concatenates its name,
molecular formula (for compounds), EC number (for enzymes), cross-reference
values, and free-text description. Reactions are excluded from the vector
index because reaction identity is better captured by stoichiometric
connectivity, which is available through the SQLite layer. Compounds, enzymes,
and pathways are indexed.

The \texttt{MetaIndex.search(query, k)} method embeds the query string with
the same model and returns the \texttt{k} approximate nearest neighbours,
along with their cosine distances. Results are then joined against SQLite to
return full node metadata.

Users may substitute any encoder that implements the \texttt{Embedder}
abstract class, which requires only a single method:

\begin{lstlisting}[style=python]
class Embedder(Protocol):
    def encode(self,
               texts: list[str]
               ) -> list[list[float]]: ...
\end{lstlisting}

% ---------------------------------------------------------------------------
\section{Query API}
\label{sec:query}
% ---------------------------------------------------------------------------

\textsc{MetaKG} exposes four query modalities through a unified API. These
operations share the same underlying graph but use different access patterns
suited to different analysis tasks. Together, they support both precision
(exact structural queries) and exploration (semantic discovery).

\subsection{Node Retrieval and Neighbourhood Traversal}

\texttt{MetaStore.node(id)} fetches a single node by internal ID. The
returned dictionary mirrors the \texttt{meta\_nodes} schema, with the
\texttt{xrefs} and \texttt{stoichiometry} fields pre-decoded from JSON.

\texttt{MetaStore.neighbours(id, rels=...)} returns all nodes reachable
from the given node by following the specified relation types. The default
relation tuple is \texttt{(SUBSTRATE\_OF, PRODUCT\_OF, CATALYZES, CONTAINS)}.
A single SQL query over the \texttt{meta\_edges} table resolves both
outgoing and incoming edges and joins the results against \texttt{meta\_nodes}.

\subsection{Reaction Detail}

\texttt{MetaStore.reaction\_detail(id)} assembles a structured view of a
single reaction, returning a dictionary with keys \texttt{substrates},
\texttt{products}, \texttt{enzymes}, \texttt{inhibitors}, \texttt{activators},
and \texttt{pathways}. Each value is a list of node dictionaries. This serves
as the primary access point for stoichiometric analysis.

\subsection{Shortest-Path Search}

\texttt{MetaStore.find\_shortest\_path(a, b, max\_hops)} performs an
iterative breadth-first search over the SQLite graph, alternating between
\texttt{SUBSTRATE\_OF} and \texttt{PRODUCT\_OF} edges to traverse the
bipartite compound--reaction graph. The search terminates when it reaches the
target node or exceeds the hop limit. Both internal IDs and external
identifiers (resolved through \texttt{xref\_index}) are accepted as
arguments.

The algorithm runs entirely in Python with SQL queries per BFS frontier,
which is practical for graphs typical of a single organism's curated metabolic
network (tens of thousands of nodes). For corpora with hundreds of thousands
of reactions, a specialised graph engine would be more appropriate.

\subsection{Semantic Search}

\texttt{MetaKG.query\_pathway(name, k)} performs a vector search over the
LanceDB index, filtering results to the \texttt{pathway} node kind. The raw
\texttt{MetaIndex.search(query, k)} method returns hits from all indexed
kinds. Both methods accept free-text queries in natural language; the
sentence-transformer model handles biological terminology effectively because
its training corpus includes scientific text.

\texttt{MetaKG.resolve\_id(s)} provides unified look-up that accepts any of:
an internal ID (\texttt{cpd:kegg:C00022}), a shorthand external reference
(\texttt{kegg:C00022}), or a display name (\texttt{Pyruvate}). It queries
\texttt{xref\_index} first for exact matches, then falls back to a
case-insensitive name search in \texttt{meta\_nodes}. This lets the same
user-facing functions accept identifiers from any source database.

\subsection{Metabolic Simulations}

Building on the structural query layer, \textsc{MetaKG} provides three
simulation modalities:

\begin{enumerate}[leftmargin=*,itemsep=2pt]
  \item \textbf{Flux Balance Analysis (FBA)} ---
    \texttt{MetaKG.simulate\_fba(pathway\_id, maximize=True)} performs
    steady-state optimisation using the stoichiometry stored in the graph.
    The result includes a status flag, objective value, and per-reaction flux
    distribution.

  \item \textbf{Kinetic ODE Integration} ---
    \texttt{MetaKG.simulate\_ode(pathway\_id, t\_end, t\_points, ...)}
    runs time-course simulation with Michaelis-Menten rate laws and seeded
    kinetic parameters (Km, Vmax, kcat). It uses an implicit stiff solver
    (BDF) tuned for metabolic systems and returns time-course concentration
    trajectories for all compounds.

  \item \textbf{What-If Perturbation Analysis} ---
    \texttt{MetaKG.simulate\_whatif(pathway\_id, scenario\_json, mode)}
    compares a baseline pathway against a perturbed version with enzyme
    knockouts, inhibition, or substrate concentration overrides. Both FBA
    and ODE modes are supported.
\end{enumerate}

Kinetic parameters are populated on first use by seeding from literature
sources (BRENDA, SABIO-RK, published metabolic models) via
\texttt{MetaKG.seed\_kinetics()}.

% ---------------------------------------------------------------------------
\section{Visualisation}
\label{sec:viz}
% ---------------------------------------------------------------------------

\subsection{2D Web Explorer}

The \texttt{metakg-viz} command launches a Streamlit~\citep{streamlit2019}
web application with three views:

\begin{enumerate}[leftmargin=*,itemsep=2pt]
  \item \textbf{Graph Browser} --- an interactive network rendered with
    pyvis~\citep{pyvis2021} and displayed in the browser. Nodes are
    colour-coded by kind (pathway: blue, reaction: red, compound: green,
    enzyme: orange) and edges by relation type. A sidebar provides filters
    by node kind and limits the number of rendered nodes to keep the
    visualisation manageable for large graphs.
  \item \textbf{Semantic Search} --- a free-text query box that calls
    \texttt{MetaIndex.search} and displays ranked results with similarity
    scores.
  \item \textbf{Node Details} --- clicking any node opens a detail panel
    showing all metadata and the node's immediate neighbourhood.
\end{enumerate}

\subsection{3D Visualiser}

The \texttt{metakg-viz3d} command launches a PyVista~\citep{sullivan2019pyvista}
interactive 3D viewer. Two layout strategies are available:

\textbf{Allium layout.} Each pathway node sits on a flat Fibonacci annulus in
the XY-plane~\citep{vogel1979fibonacci}. Reaction and compound nodes belonging
to a pathway are distributed on a Fibonacci sphere centred on the pathway's
position, creating a visual metaphor of an inflorescence. Nodes belonging to
multiple pathways are placed at the centroid of their pathway positions, so
cross-pathway metabolites appear in intermediate locations.

\textbf{LayerCake layout.} Nodes are stratified by kind along the Z-axis:
pathways occupy the lowest layer, reactions the middle, and compounds and
enzymes the top. Within each layer, nodes follow a golden-angle spiral to
minimise overlap. This layout suits inspection of the bipartite
compound--reaction structure.

Both layouts export to HTML (for web reports) and PNG (for publication
figures).

% ---------------------------------------------------------------------------
\section{Model Context Protocol Server}
\label{sec:mcp}
% ---------------------------------------------------------------------------

The Model Context Protocol (MCP)~\citep{anthropic2024mcp} is a lightweight
JSON-RPC standard that lets large-language-model assistants call typed tool
functions. \textsc{MetaKG} implements an MCP server exposing the knowledge
graph through nine typed tools:

\begin{description}[leftmargin=*,itemsep=2pt]
  \item[\texttt{query\_pathway(name, k)}] Semantic pathway search. Returns
    up to \texttt{k} pathway nodes whose descriptions are closest to the
    query in embedding space, along with their member-reaction counts.
  \item[\texttt{get\_compound(id)}] Returns a compound node with its
    connected reactions, accepting any supported identifier format
    (full URI, shorthand, or display name).
  \item[\texttt{get\_reaction(id)}] Returns full stoichiometric detail for
    a single reaction, including substrates, products, and catalysing enzymes.
  \item[\texttt{find\_path(a, b, max\_hops)}] Returns the shortest metabolic
    path between two compounds via bidirectional BFS.
  \item[\texttt{simulate\_fba(pathway\_id, ...)}] Runs Flux Balance Analysis
    on a pathway and returns the optimal flux distribution and shadow prices.
  \item[\texttt{simulate\_ode(pathway\_id, ...)}] Runs kinetic ODE simulation
    (BDF solver) and returns time-course concentration trajectories.
  \item[\texttt{simulate\_whatif(pathway\_id, scenario\_json, mode)}]
    Compares baseline and perturbed simulations; supports enzyme knockouts,
    activity scaling, and substrate overrides.
  \item[\texttt{get\_kinetic\_params(reaction\_id)}] Retrieves stored kinetic
    and regulatory parameters (Km, Vmax, kcat, Ki, $\Delta G^{\circ\prime}$) for a reaction.
  \item[\texttt{seed\_kinetics(force)}] Populates the database with curated
    literature kinetic parameters from BRENDA, SABIO-RK, and published models.
\end{description}

The server starts with:

\begin{lstlisting}[style=bash]
metakg-mcp --db .metakg/meta.sqlite \
           --transport stdio
\end{lstlisting}

\noindent and communicates over standard input/output (the \texttt{stdio}
transport) or as an HTTP server-sent events stream (the \texttt{sse}
transport). The \texttt{stdio} transport is the standard configuration for
Claude~\citep{anthropic2024claude} and compatible MCP clients.

% ---------------------------------------------------------------------------
\section{Worked Example: Complete Human Metabolome}
\label{sec:usage}
% ---------------------------------------------------------------------------

We demonstrate \textsc{MetaKG} on the complete \textit{Homo sapiens}
metabolome: all 369 KEGG pathways (metabolic, signaling, and regulatory).
This corpus spans central carbon metabolism (glycolysis/gluconeogenesis, TCA
cycle, pentose phosphate pathway, fatty acid degradation, oxidative
phosphorylation), amino acid and nucleotide metabolism, secondary metabolism,
carbohydrate metabolism, lipid metabolism, and signaling networks.

\subsection{Building the Knowledge Graph}

\begin{lstlisting}[style=bash, caption={Building the complete human metabolome knowledge graph.}]
$ metakg-build --data ./data/hsa_pathways
Building MetaKG from ./data/hsa_pathways...
data_root   : ./data/hsa_pathways
db_path     : .metakg/meta.sqlite
nodes       : 22290  {'compound': 5115, 'enzyme': 14667,
              'pathway': 369, 'reaction': 2139}
edges       : 11298  {'CATALYZES': 2406, 'CONTAINS': 3809,
              'PRODUCT_OF': 2532, 'SUBSTRATE_OF': 2551}
indexed     : 20151 vectors  dim=384
\end{lstlisting}

By default, \texttt{metakg-build} wipes the existing database and vector index
before parsing (the safest, most predictable behaviour). To incrementally merge
new pathway files instead, use the \texttt{--no-wipe} flag or the convenience
command \texttt{metakg-update --data ./pathways}.

\subsection{Structural Queries via the Python API}

\begin{lstlisting}[style=python, caption={Compound retrieval and neighbourhood traversal.}]
from metakg import MetaKG

kg = MetaKG()

# Retrieve pyruvate and its connected reactions
cpd = kg.get_compound("cpd:kegg:C00022")
print(cpd["name"])
for rxn in cpd["reactions"]:
    print(f"  {rxn['name']:30s} {rxn['role']}")

# Full reaction detail
rxn = kg.get_reaction("rxn:kegg:R00200")
print(f"Substrates: {[s['name'] for s in rxn['substrates']]}")
print(f"Products:   {[p['name'] for p in rxn['products']]}")
\end{lstlisting}

Output:
\begin{lstlisting}[style=bash, caption={}]
Pyruvate
  R00703                         SUBSTRATE_OF
  R00014                         SUBSTRATE_OF
  R00431                         SUBSTRATE_OF
  R00209                         SUBSTRATE_OF
  R00200                         PRODUCT_OF
  ... and 5 more
Substrates: ['Phosphoenolpyruvate', 'ADP']
Products:   ['Pyruvate', 'ATP']
\end{lstlisting}

\subsection{Shortest-Path Search}

\begin{lstlisting}[style=python, caption={Finding the shortest metabolic route between two compounds.}]
from metakg import MetaKG

kg = MetaKG()

result = kg.find_path(
    "cpd:kegg:C00031",   # D-Glucose
    "cpd:kegg:C00022",   # Pyruvate
    max_hops=12,
)
print(f"Path length: {result['hops']} steps")
for node in result["path"]:
    print(f"  {node['kind']:10s} {node['name']}")
\end{lstlisting}

Output:
\begin{lstlisting}[style=bash, caption={}]
Path length: 9 steps
  compound   C00031
  reaction   R00299
  compound   C00092
  reaction   R00771
  compound   C00085
  reaction   R00756
  compound   C00354
  reaction   R01068
  compound   C00118
  reaction   R01061
  compound   C00236
  reaction   R01662
  compound   C01159
  reaction   R09532
  compound   C00631
  reaction   R00658
  compound   C00074
  reaction   R00200
  compound   C00022
\end{lstlisting}

The query resolves in milliseconds on the local SQLite index. The algorithm
scales efficiently to the complete human metabolome (17,050 nodes, 11,286
edges) using bidirectional BFS with early termination. Typical shortest-path
queries complete in 10--50~ms.

\subsection{Semantic Search}

\begin{lstlisting}[style=python, caption={Semantic pathway retrieval.}]
from metakg import MetaKG

kg = MetaKG()

result = kg.query_pathway("fatty acid beta-oxidation", k=5)
for hit in result.hits:
    print(f"{hit['name']:40s}  dist={hit['_distance']:.3f}")
\end{lstlisting}

Output:
\begin{lstlisting}[style=bash, caption={}]
Fatty acid degradation                    dist=1.174
alpha-Linolenic acid metabolism           dist=1.183
Fatty acid metabolism                     dist=1.200
Biosynthesis of unsaturated fatty acids   dist=1.245
Linoleic acid metabolism                  dist=1.252
\end{lstlisting}

The semantic search correctly identifies related pathways despite differences
in nomenclature between the query and the KEGG pathway names. The
sentence-transformer model handles synonymy and domain terminology effectively.

\subsection{Pathway Analysis Report}

The \texttt{metakg-analyze} command runs a seven-phase analysis and produces a
structured Markdown report:

\begin{lstlisting}[style=bash]
$ metakg-analyze --output analysis.md --top 10
\end{lstlisting}

The report covers: (1)~aggregate graph statistics; (2)~hub metabolites ranked
by degree; (3)~reactions ranked by stoichiometric complexity;
(4)~cross-pathway hub detection; (5)~pairwise pathway coupling by shared
metabolites; (6)~topological features (dead-end compounds, isolated nodes);
and (7)~enzymes ranked by reaction count. On the complete human metabolome
(369 pathways), ATP, NAD$^{+}$, coenzyme~A, and pyruvate emerge as the top
hub metabolites by degree---consistent with their known roles as central
energy and carbon carriers. The analysis also reveals cross-pathway metabolite
connectivity and identifies the reactions with highest stoichiometric
complexity.

\subsection{Mixed-Format Ingestion}

To illustrate multi-format ingestion, one can combine KGML files with a
BioPAX export from Reactome and a custom CSV table in a single data directory:

\begin{lstlisting}[style=bash]
$ ls pathways/
hsa00010.xml    # KGML (KEGG)
R-HSA-70171.owl # BioPAX (Reactome)
custom_rxns.csv # CSV (in-house data)

$ metakg-build --data ./pathways
\end{lstlisting}

The parser dispatcher examines the root XML element of each \texttt{.xml} or
\texttt{.owl} file and selects the appropriate parser; \texttt{.csv} and
\texttt{.tsv} files go to the tabular parser. Nodes sharing a KEGG or ChEBI
cross-reference across files are automatically merged in SQLite through the
xref index.

% ---------------------------------------------------------------------------
\section{Implementation Notes}
\label{sec:impl}
% ---------------------------------------------------------------------------

\textbf{Dependencies.} The core package requires Python~3.12,
\texttt{lancedb$\,\geq\,$0.29}, \texttt{numpy$\,\geq\,$1.24}, and
\texttt{sentence-transformers$\,\geq\,$2.7}. No network connection is needed
at run time once the embedding model has been downloaded. Simulation support
(FBA and ODE) requires the optional \texttt{scipy$\,\geq\,$1.11} package
installed via \texttt{--extras simulate}. BioPAX support requires the optional
\texttt{rdflib} package; the 2D and 3D visualisers need optional extras
installed via \texttt{poetry install --extras viz} or \texttt{--extras viz3d}.

\textbf{Installation.}

\begin{lstlisting}[style=bash]
git clone https://github.com/Flux-Frontiers/meta_kg
cd meta_kg
poetry install --extras all
\end{lstlisting}

\textbf{Thread safety.} The SQLite connection uses WAL mode with a single
Python object per process. The current implementation is not thread-safe;
callers should create one \texttt{MetaKG} instance per process or protect a
shared instance with a lock.

\textbf{Incremental rebuild.} The stable ID scheme makes incremental builds
straightforward. Running \texttt{metakg-build} with \texttt{--no-wipe} (or using
the convenience command \texttt{metakg-update}) issues \texttt{INSERT OR REPLACE}
statements, updating existing nodes and appending new ones without duplicating entries.

\textbf{Codebase self-analysis.} \textsc{MetaKG} is itself analysed with its
sister tool \textsc{CodeKG}~\citep{fluxfrontiers2025codekg}, which builds a
structural and semantic knowledge graph of the Python source code. This
enables navigating the MetaKG implementation via natural-language queries and
confirms that the two tools share compatible architectural patterns. On the
\textsc{MetaKG} codebase, the \textsc{CodeKG} static analysis produces 3,136
nodes and 2,920 edges spanning 27 modules, embedded into a 384-dimensional
vector index of 290 vectors.

% ---------------------------------------------------------------------------
\section{Discussion}
\label{sec:discussion}
% ---------------------------------------------------------------------------

\subsection{Architectural Design Choices and Novel Aspects}

Several deliberate design trade-offs distinguish \textsc{MetaKG} from existing
systems. First, the dual-layer architecture (SQLite for structure, LanceDB for
semantics) is novel in the metabolic pathway domain. Graph databases like
Neo4j handle complex queries but carry operational overhead (server
management, scaling, deployment) and do not address multi-format parsing.
Specialised vector databases like Weaviate or Pinecone excel at semantic
search but are poor fits for structural graph traversal and require network
access. By pairing lightweight SQLite with local vector search,
\textsc{MetaKG} delivers both capabilities in a self-contained package suited
to exploratory research and reproducible analysis.

Second, the deterministic identifier scheme with synthetic hashing enables
reproducible cross-format merging without a centralised reconciliation
service. Unlike MetaNetX (which requires API calls to reconcile identifiers),
\textsc{MetaKG} produces self-contained graphs. The same input files always
yield the same output graph, making the knowledge graph a version-controlled
artefact that supports reproducible science and offline workflows.

Third, the four-modality query interface (structural, pathfinding, semantic,
stoichiometric) is intentionally broad. An analyst can start with a
natural-language semantic query (``fatty-acid beta-oxidation''), then drill
into structural detail (shortest paths, stoichiometric coefficients) without
switching tools. This contrasts with systems that specialise in a single query
paradigm.

Fourth, the Model~Context~Protocol integration is forward-looking. As
large-language-model assistants become standard tools in computational biology,
making the knowledge graph a first-class data source for Claude, ChatGPT, and
future assistants is a natural step. The MCP interface represents a design
commitment to AI-accessible knowledge graphs, not merely another API layer.

\subsection{Scope and Design Trade-offs}

\textbf{Snapshot-based operation.} \textsc{MetaKG} is a local analysis tool,
not a live database mirror. The knowledge graph is a snapshot at build time;
users must re-run \texttt{metakg-build} after updating source files. This
keeps the system self-contained and free of dependencies on external services
during analysis. For research workflows where reproducibility matters most,
this is an advantage. For applications requiring real-time data updates, a
live database backend would be more suitable.

\textbf{Scale.} SQLite and in-process BFS handle graphs of up to roughly
100,000 nodes, covering full reconstructed metabolic networks for a single
organism. For pan-genome or multi-species analyses---where node counts reach
into the millions---a dedicated graph database engine (e.g.,
Neo4j~\citep{neo4j2020} or K\`uzu~\citep{feng2023kuzu}) and a distributed
vector index would be preferable. The storage layer is designed to be
replaceable: \texttt{MetaStore} and \texttt{MetaIndex} are concrete classes
behind well-defined interfaces, so alternative backends can be substituted
without changing the parser or query layers.

\textbf{Identifier reconciliation.} The current cross-reference merge relies
on exact match of external identifiers. Two compounds sharing a biological
identity but differing in stereochemistry or protonation state will not merge
automatically. More thorough reconciliation would require integration with a
name normalisation service such as
MetaNetX~\citep{moretti2021metanetx} or UMLS~\citep{bodenreider2004umls}.

\textbf{Stoichiometric models and simulations.} \textsc{MetaKG} provides
three simulation modalities: (1)~Flux Balance Analysis (FBA) for steady-state
flux optimisation; (2)~kinetic ODE integration with Michaelis-Menten rate laws
using an implicit stiff solver (BDF) tuned for metabolic pathways; and
(3)~what-if perturbation analysis for enzyme knockouts, inhibition, and
substrate overrides. Kinetic parameters are seeded from literature sources
(BRENDA, SABIO-RK, published models). For advanced constraint-based modelling,
the SBML parser preserves all information needed to reconstruct a COBRApy
model.

\textbf{Performance on complete human metabolome.} The complete
\textit{Homo sapiens} metabolome (369 KEGG pathways, 17,050 nodes, 11,286
edges, 20,151 vector embeddings) shows the following performance
characteristics:

\begin{table}[H]
  \centering
  \caption{Performance on the complete human metabolome (17K nodes, 11K edges).
  All timings measured on a 2024 MacBook Pro (Apple M3 Max, 36\,GB RAM, 1\,TB SSD).}
  \label{tab:performance}
  \small
  \begin{tabular}{@{}lll@{}}
    \toprule
    Operation & Time & Details \\
    \midrule
    Build graph (parse + index) & 30--60s & All 369 pathways, LanceDB + SQLite \\
    Semantic search (natural language) & 100--500ms & Vector similarity on 20K nodes \\
    Shortest-path (6 hops) & 10--50ms & BFS on 11K edges \\
    ODE simulation (10 units) & 150--400ms & BDF solver, 24 compounds \\
    Streamlit rerun & 0.5--1.5s & Batch query + session cache \\
    \bottomrule
  \end{tabular}
\end{table}

All operations complete within practical timeframes for interactive
exploration. The graph is well suited to research workflows demanding
reproducible analysis, offline access, and programmatic querying.

\textbf{Future directions.} Planned extensions include graph-theoretic
centrality measures (betweenness, closeness, eigenvector centrality) via
NetworkX for hub ranking; a GraphQL query endpoint; integration with the
UniProt and ChEBI REST APIs for on-demand annotation; differential pathway
analysis across organisms; GPU-accelerated embedding for large datasets; and
export to Cytoscape~\citep{shannon2003cytoscape} JSON for interoperability
with the broader network biology ecosystem. Advanced kinetic parameter
optimisation (fitting from experimental data) is also planned.

% ---------------------------------------------------------------------------
\section{Conclusion}
\label{sec:conclusion}
% ---------------------------------------------------------------------------

\textsc{MetaKG} fills a gap in metabolic data integration. Existing systems
force a choice among convenient web interfaces (limited programmability, no
semantic search, web-dependent), specialised graph databases (high operational
complexity, parsing unsolved), and reconciliation services (no queryable graph
provided). \textsc{MetaKG} breaks this pattern with a dual-layer local
knowledge graph that unifies multi-format pathway data and delivers four
orthogonal query modalities---structural, pathfinding, semantic, and
stoichiometric---through a single API, CLI, and LLM-accessible MCP interface.

The core contributions are: (1)~a stable, deterministic identifier scheme
enabling reproducible cross-format merging without external services; (2)~a
dual-layer storage architecture (SQLite + LanceDB) that eliminates the forced
choice between relational precision and semantic expressivity; (3)~a
two-phase name-enrichment pipeline that replaces bare KEGG accessions with
human-readable names, making all query results immediately interpretable
without network access; (4)~four unified query modalities accessible to
analysts at all levels of programming expertise; and (5)~a forward-looking
MCP server exposing nine typed tools that make metabolic knowledge graphs
first-class data sources for AI assistants.

The design philosophy---local-first, self-contained, snapshot-based---is
deliberate and represents a clean separation from live database mirrors. This
makes \textsc{MetaKG} especially well suited to research workflows where
reproducibility, offline analysis, and version control matter most. For
applications needing real-time data, larger graph corpora, or distributed
deployment, the modular architecture allows the storage and index backends to
be swapped.

We expect \textsc{MetaKG} to be immediately useful as: (1)~a foundation for
pathway analysis scripts; (2)~a data preparation stage for machine-learning
workflows; (3)~an AI-accessible knowledge source for metabolic reasoning in
large-language-model applications; and (4)~a template for similar knowledge
graph systems in related biological domains (protein interactions, gene
regulatory networks, drug--target interactions).

The software is freely available at
\url{https://github.com/Flux-Frontiers/meta_kg} under the Elastic License~2.0.

% ---------------------------------------------------------------------------
% Appendix: Format-Specific Parser Details
% ---------------------------------------------------------------------------
\appendix

\section{Format-Specific Parsers}
\label{sec:parsers-detail}

All parsers conform to an abstract base class:

\begin{lstlisting}[style=python]
class PathwayParser:
    def can_handle(self, path: Path) -> bool: ...
    def parse(self, path: Path
              ) -> tuple[list[MetaNode],
                         list[MetaEdge]]: ...
\end{lstlisting}

Parsers are stateless and pure: the same input file always produces the same
output. The \texttt{MetabolicGraph} layer caches the combined node and edge
lists after the first parse, so repeated calls do not re-read disk.

\subsection{KGML Parser}

KEGG Markup Language files are the native export format of the KEGG pathway
database~\citep{kanehisa2021kegg}. Each file is an XML document whose root
element is \texttt{<pathway>}. The parser uses the Python standard-library
\texttt{ElementTree} module (no third-party XML dependency) and extracts three
kinds of child elements:

\begin{itemize}[leftmargin=*,itemsep=2pt]
  \item \texttt{<entry>} elements with \texttt{type="compound"} become
    \texttt{compound} nodes.
  \item \texttt{<entry>} elements with \texttt{type="gene"} or
    \texttt{type="enzyme"} become \texttt{enzyme} nodes.
  \item \texttt{<reaction>} elements become \texttt{reaction} nodes with
    their \texttt{<substrate>} and \texttt{<product>} children encoded as
    stoichiometry JSON. The enclosing pathway becomes a \texttt{CONTAINS}
    edge to each reaction.
\end{itemize}

Format detection checks the root element tag rather than the file extension,
making the parser robust to KEGG files served without the \texttt{.kgml}
extension.

\subsection{SBML Parser}

The Systems Biology Markup Language~\citep{keating2020sbml} is the standard
serialisation format for constraint-based metabolic models built with tools
such as COBRApy~\citep{ebrahim2013cobrapy}. SBML Level~2 and~3 files share a
common XML namespace ending in \texttt{sbml}; the parser detects the format by
matching the root element's local name.

Species elements map to \texttt{compound} nodes. Reaction elements map to
\texttt{reaction} nodes. Stoichiometry comes from
\texttt{<listOfReactants>} and \texttt{<listOfProducts>} children. Modifier
species are classified by their SBO~term~\citep{courtot2011sbo}: SBO:0000013
(catalyst) generates \texttt{CATALYZES} edges; SBO:0000020 (inhibitor)
generates \texttt{INHIBITS} edges; other modifiers generate
\texttt{ACTIVATES} edges.

\subsection{BioPAX Parser}

Biological Pathway Exchange Level~3~\citep{demir2010biopax} is an OWL
ontology serialised as RDF/XML, used by Reactome~\citep{jassal2020reactome},
WikiPathways~\citep{martens2021wikipathways}, and the NCI Pathway Interaction
Database. Parsing requires the optional \texttt{rdflib} dependency, installed
via the \texttt{biopax} extra. The parser performs SPARQL-style pattern
matching over the RDF graph to extract:

\begin{itemize}[leftmargin=*,itemsep=2pt]
  \item \texttt{SmallMolecule} instances $\to$ \texttt{compound} nodes.
  \item \texttt{Protein} instances $\to$ \texttt{enzyme} nodes.
  \item \texttt{BiochemicalReaction} instances $\to$ \texttt{reaction} nodes,
    with \texttt{left}/\texttt{right} properties becoming substrate and
    product edges, and \texttt{controller} properties becoming
    \texttt{CATALYZES} edges.
  \item \texttt{Pathway} instances $\to$ \texttt{pathway} nodes, with
    \texttt{memberPathwayComponent} links becoming \texttt{CONTAINS} edges.
\end{itemize}

\subsection{CSV/TSV Parser}

For custom or unpublished data, \textsc{MetaKG} accepts flat tables with a
configurable column schema. The default column layout is:

\begin{center}
\small
\texttt{reaction\_id}, \texttt{reaction\_name}, \texttt{substrate},
\texttt{product}, \texttt{enzyme}, \texttt{stoich\_substrate},
\texttt{stoich\_product}, \texttt{pathway}, \texttt{ec\_number},
\texttt{substrate\_formula}, \texttt{enzyme\_uniprot}
\end{center}

Multiple rows with the same \texttt{reaction\_id} merge into a single
reaction node---the standard way to encode multi-substrate or multi-product
reactions in tabular form. A \texttt{CSVParserConfig} dataclass allows
remapping all column names, making the parser suitable for lab-produced
spreadsheets and bulk downloads from custom databases.

% ---------------------------------------------------------------------------
% Bibliography
% ---------------------------------------------------------------------------
\bibliographystyle{plainnat}
\bibliography{references}

\end{document}
